\section{Are goal heads anything special? How many do we need for a decent readout?}
\label{app:readout_sources}
\label{sec:heads:residual}
\arnab{baseline comparison against full attention module outputs and residual. all with additive\_norm and $\alpha=0.1$ to ensure ~fairness.}
\arnab{readouts with different number of heads, across different pivot positions. the figure on the main text should be the aggregated numbers across the 4 pivot points.}

\begin{figure}[ht]\centering
\includegraphics[width=\linewidth]{Figures/appendix/source_comparison.pdf}
\caption{\textbf{Comparison of injected sources.} Norm-relative injection at $\alpha=.1$ on the locked corpus, pooled over five positions ($n=1,265$ per bar). Layer endpoints are inclusive: the saved half-open source ranges \texttt{[40,50)} and \texttt{[25,56)} correspond to L40--49 and L25--55. The attention and residual sources are means over their stated layers. Poor raw-residual injection is a result about this probe configuration, not proof that the residual lacks subgoal information.}
\label{tab:heads-vs-residual}
\label{fig:heads-vs-residual}
\end{figure}
Selected contributions provide a substantially more effective injection direction than the tested full-attention and raw-residual aggregates. Equal relative norm controls overall perturbation magnitude, but does not ensure that a shared injection strength or aggregation rule is optimal for every source. OVReader provides a separate test of information extracted from the residual.

\clearpage
\subsection{}
\label{app:sparsity_pivots}
\begin{figure}[!ht]
\centering
\includegraphics[width=\linewidth]{Figures/sparsity_other_pivots.pdf}
\caption{Head-budget curves for the remaining four capture positions, using the same 253 specifications and configurations as \Cref{fig:sparsity}. Stars indicate the 122-head reference; blue solid lines show readout faithfulness and orange dashed lines implementation transfer. The first implementation token is a comparison position, not evidence of encoding before implementation. The common vertical scale makes differences across positions visible.}
\label{fig:sparsity-other-pivots}
\end{figure}
We recompute each rate as the fraction of saved oracle verdicts equal to \texttt{YES}, including all 253 records per position in its denominator. The same specification identifiers occur at every head budget in both evaluation modes. Each position is a repeated measurement of those specifications. All runs use norm-relative additive injection at $\alpha=.1$, layers 35--45, and steering throughout probe decoding. The head selections differ in construction and are not strictly nested; the 122-head reference is displayed separately from the connected sweep.



\section{Across Model Comparisons}

\subsection{Model Size}
\arnab{compare Qwen-2.5-Coder-instruct, 14B vs 32B}

\subsection{Code-finetuning helps}
\arnab{compare Qwen-2.5 Coder vs non-coder. both 32B}

\subsection{RL-finetuning also helps}
\arnab{compare OLMo models. RL-tuned vs SFT. Causality increases a bit with RL but faithfulness goes down. RL makes OLMo almost always generate its response as \textit{"I cannot do it ..."}, a table with a 3/4 representitive readouts.}